\SetPageTheme{story}
\PageHeader
  {Calculatorul nu știe să numere}
  {Când jocul pornește, calculatorul nu vede eroi, uși sau hărți. Vede doar stări pe care trebuie să le țină minte.}
  {Arc narativ 1 din 3}

\begin{missionbox}[Intrarea în poveste]
Jucătorul apare pe hartă. Are vieți, are scor, are o poziție, are un inventar. Toate aceste lucruri trebuie ținute în memorie, altfel jocul uită ce s-a întâmplat după fiecare cadru (frame).

Primul pas de programator nu este să deseneze efecte speciale. Primul pas este să fixeze starea inițială.
\end{missionbox}

\begin{conceptbox}[Reținem!]
\begin{itemize}
\item O \textbf{variabilă} reține o stare curentă a jocului.
\item Fără stare, nu există progres, memorie sau logică.
\item Schimbarea vizibilă de pe ecran înseamnă, de fapt, o schimbare invizibilă în memorie.
\end{itemize}
\end{conceptbox}

\BlocklyExperiment{cq-i1-ram-cells}{Memoria la pornire}{Leagă blocurile de inițializare cu starea jocului: vieți, scor, nivel, chei. Observă cum fiecare bloc „rezervă” un loc în memoria calculatorului.}

\clearpage

\SetPageTheme{story}
\PageHeader
  {Cum gândește mașina}
  {Mașina nu judecă dacă o alegere este bună sau rea. Ea doar citește valori și execută reguli matematice.}
  {Arc narativ 2 din 3}

\begin{machinebox}[Model intern simplificat]
\begin{lstlisting}[style=codequest]
vieti <- 3
scor <- 0
nivel <- 1
chei <- 0
\end{lstlisting}
\end{machinebox}

\begin{missionbox}[Ce înseamnă asta în joc]
Când personajul pierde o viață, el nu „cade” de pe ecran ca o idee abstractă. În memorie, variabila \texttt{vieti} scade cu o unitate. Când găsește un cristal rar, variabila \texttt{scor} crește.

Această legătură directă dintre scena vizibilă și starea invizibilă este firul narativ care ne va însoți pe parcursul întregii reviste.
\end{missionbox}

\begin{guidebox}[Primul check]
Pornește de la valorile de mai sus și urmărește logica:
\begin{enumerate}
\item După o greșeală și o coliziune: ce valoare va avea \texttt{vieti}?
\item După colectarea a două cristale: ce valoare va avea \texttt{scor}?
\item După trecerea cu succes a primului nivel: ce devine \texttt{nivel}?
\end{enumerate}
\AnswerSpace{4}
\end{guidebox}

\clearpage

\SetPageTheme{story}
\PageHeader
  {De la poveste la control}
  {Când starea este clară, putem controla întregul traseu al jocului: urcare în nivel, blocaje, restart sau bonusuri.}
  {Arc narativ 3 din 3}

\begin{checkpointbox}[Legătura care rămâne]
Până la finalul acestei ediții, fiecare idee nouă de programare va răspunde la aceeași întrebare fundamentală:
\begin{center}
\textbf{Ce variabilă se schimbă și de ce?}
\end{center}
\end{checkpointbox}

\begin{semibox}[Exercițiu de tranziție]
Imaginează-ți un mini-joc cu un labirint întunecat. Scrie trei stări pe care ai vrea să le urmărești:
\begin{enumerate}
\item o stare care măsoară \textbf{progresul} (cât de departe ai ajuns);
\item o stare care măsoară \textbf{riscul} (cât de aproape este pericolul);
\item o stare care măsoară \textbf{resursele} (ce ai adunat pentru a supraviețui).
\end{enumerate}
\AnswerSpace{4}
\end{semibox}

\SimulatorCard{cq-i1-ram-cells}{RAM și celule de memorie}{Cum fiecare variabilă ocupă o „celulă” de memorie, iar eticheta (numele) rămâne aceeași chiar dacă valoarea din interior se modifică constant.}{Mișcă sliderul pentru \texttt{vieti} și \texttt{scor} în simulator și observă cum celulele de memorie își schimbă culoarea și conținutul. Poți descrie scena din joc doar citind aceste numere?}

\clearpage

\SetPageTheme{theory}
\PageHeader
  {Variabile și memorie}
  {O variabilă bine aleasă spune limpede ce ține minte jocul și de ce informația aceea contează.}
  {Variabile 1 din 5}

\begin{conceptbox}[Definiție utilă]
O variabilă are trei componente esențiale pe care trebuie să le stăpânești:
\begin{itemize}
\item \textbf{Numele} (cum o identifici în cod, de exemplu \texttt{viteza});
\item \textbf{Valoarea curentă} (ce conține acum, de exemplu \texttt{10});
\item \textbf{Rolul} (ce poveste spune în logica jocului, de exemplu „cât de repede se mișcă inamicul”).
\end{itemize}
\end{conceptbox}

\begin{examplebox}[Exemplu minim de declarare]
\begin{lstlisting}[style=codequest]
scor <- 0
vieti <- 3
energie <- 100
\end{lstlisting}
\end{examplebox}

\SyntaxCheck{variabile}

\clearpage

\SetPageTheme{guided}
\PageHeader
  {Citim memoria corect}
  {Ne interesează întotdeauna valoarea de acum, nu istoricul complet al tuturor valorilor care au fost acolo înainte.}
  {Variabile 2 din 5}

\begin{guidebox}[Ghidat: Suprascrierea]
Pornim cu următoarea stare inițială:
\begin{lstlisting}[style=codequest]
monede <- 4
chei <- 1
\end{lstlisting}
După un eveniment de tip „colectare”, se execută:
\begin{lstlisting}[style=codequest]
monede <- 9
chei <- 2
\end{lstlisting}
\begin{enumerate}
\item Ce valori rămân în memorie la finalul execuției?
\item Care valoare veche a fost „ștearsă” sau înlocuită?
\item De ce crezi că nu putem avea două valori simultan pentru aceeași variabilă?
\end{enumerate}
\AnswerSpace{6}
\end{guidebox}

\begin{guidebox}[Ghidat cu desen: Tabel de stare]
Completează tabelul de mai jos urmărind valorile finale dintr-o sesiune de joc:
\begin{center}
\begin{tabular}{lll}
\toprule
Variabilă & Valoare Inițială & Valoare Finală \\
\midrule
\texttt{scor} & 5 & \\
\texttt{vieti} & 3 & \\
\texttt{nivel} & 1 & \\
\bottomrule
\end{tabular}
\end{center}
\AnswerSpace{4}
\end{guidebox}

\clearpage

\SetPageTheme{simulation}
\PageHeader
  {Vedem schimbarea în timp real}
  {Când modifici un slider în simulator, vezi imediat cum se rescrie starea în „măruntaiele” calculatorului.}
  {Variabile 3 din 5}

\SimulatorCard{cq-i1-variable-state}{Evoluția unei variabile}{Observă cum aceeași variabilă își schimbă valoarea după fiecare eveniment important din joc, fără a crea variabile noi.}{Simulează trei evenimente consecutive: un bonus, o coliziune și atingerea unui checkpoint. Notează după fiecare pas valoarea nouă pe care o ia variabila în memorie.}

\begin{semibox}[Analiză după simulare]
\begin{enumerate}
\item Care dintre evenimentele simulate a produs cea mai mare schimbare numerică?
\item Ce variabilă ai urmări cu cea mai mare atenție dacă ai vrea să detectezi starea de \textit{Game Over}?
\item Ce variabilă ar indica cel mai clar că jucătorul a câștigat runda?
\end{enumerate}
\AnswerSpace{5}
\end{semibox}

\clearpage

\SetPageTheme{guided}
\PageHeader
  {Exersăm cu sprijin redus}
  {Acum este rândul tău să completezi stări inițiale și să gândești tranziții realiste pentru un joc.}
  {Variabile 4 din 5}

\begin{semibox}[Construiești un set de stări]
Construiește un set de variabile pentru un joc de tip „platformer” (precum Mario):
\begin{enumerate}
\item Alege 5 variabile relevante pentru acest gen de joc;
\item Dă-i fiecăreia o valoare inițială logică;
\item Scrie pentru fiecare variabilă un eveniment concret care i-ar schimba valoarea.
\end{enumerate}
\AnswerSpace{8}
\end{semibox}

\begin{semibox}[Urmărești schimbarea pas cu pas]
Pornești cu o variabilă numită \texttt{energie <- 6}. După fiecare pas de mai jos, notează noua valoare a energiei:
\begin{enumerate}
\item Personajul lovește un obstacol spinos;
\item Personajul colectează un scut protector;
\item Personajul folosește o abilitate specială de viteză.
\end{enumerate}
\AnswerSpace{6}
\end{semibox}

\clearpage

\SetPageTheme{challenge}
\PageHeader
  {Rezolvi independent}
  {Pagina aceasta este spațiul tău pentru o soluție completă, explicată clar, cu spațiu real de lucru.}
  {Variabile 5 din 5}

\begin{solobox}[Fără ghidaj: Creație proprie]
Imaginează-ți un mini-joc complet nou (de exemplu, unul cu avioane sau cu gătit) și definește-i structura:
\begin{enumerate}
\item Identifică minimum 6 variabile relevante;
\item Stabilește valori inițiale pentru toate;
\item Scrie o secvență de 8 linii de cod (pseudocod) în care valorile se schimbă logic;
\item Explică în 5 rânduri cum se „simte” jocul doar citind aceste schimbări.
\end{enumerate}
\AnswerSpace{12}
\end{solobox}

\begin{solobox}[Reflexia programatorului]
Care este diferența fundamentală dintre a avea „un nume de variabilă” și a avea „o variabilă bine aleasă” pentru claritatea codului tău?
\AnswerSpace{5}
\end{solobox}

\clearpage
